\documentclass[a4paper,11pt]{article}
\usepackage[utf8]{inputenc}
\usepackage[T1]{fontenc}
\usepackage[french]{babel}
\usepackage[margin=1.5cm]{geometry}
\usepackage{amsmath,amsfonts,amssymb}
\usepackage{tabularx}
\usepackage{enumitem}

% Ajustement de la hauteur des lignes du tableau pour laisser la place d'écrire
\renewcommand{\arraystretch}{2.2}

\begin{document}

\pagestyle{empty}

% ==========================================
% FICHE 4
% ==========================================
\begin{center}
    \Large\textbf{Entraînement : Automatismes 04}
\end{center}

\noindent \textbf{Partie 1 \hfill 6 points \hfill 20 minutes}

\medskip
\noindent \textbf{Consigne :} Pour chaque question, recopier sur la copie son numéro et la réponse correspondante. Aucune justification n'est demandée.

\bigskip

\noindent
\begin{tabularx}{\textwidth}{|c|X|p{4cm}|}
\hline
\textbf{N$^{\circ}$} & \textbf{Énoncé} & \textbf{Réponse} \\ \hline
\textbf{1} & Écrire sous la forme d'une seule puissance : $7^3 \times 7^5$. & \\ \hline
\textbf{2} & Développer et réduire l'expression : $A = 3(2x - 5)$. & \\ \hline
\textbf{3} & Dans un triangle $ABC$ rectangle en $A$, on sait que $AB = 3\text{ cm}$ et $AC = 4\text{ cm}$. Calculer la longueur de l'hypoténuse $BC$. & \\ \hline
\textbf{4} & On lance un dé à 6 faces non truqué. Quelle est la probabilité d'obtenir un nombre strictement supérieur à 4 ? & \\ \hline
\textbf{5} & Donner l'écriture scientifique du nombre : $45\,000$. & \\ \hline
\textbf{6} & Factoriser l'expression suivante : $B = x^2 + 6x$. & \\ \hline
\textbf{7} & Les droites $(DE)$ et $(BC)$ sont parallèles. On sait que $AD = 2\text{ cm}$, $AB = 6\text{ cm}$ et $AE = 3\text{ cm}$. Calculer $AC$. & \\ \hline
\textbf{8} & Quelle est l'image de $3$ par la fonction $f$ définie par $f(x) = x^2 - 2$ ? & \\ \hline
\textbf{9} & Un pavé droit a pour base un carré de côté $5\text{ cm}$ et pour hauteur $10\text{ cm}$. Quel est son volume ? & \\ \hline
\textbf{10} & Résoudre l'équation : $4x + 7 = 19$. & \\ \hline
\end{tabularx}

\newpage

\subsection*{Correction (pour s'auto-évaluer)}
\begin{enumerate}[label=\textbf{\arabic*.}]
    \item $\mathbf{7^8}$ (car on additionne les exposants : $3 + 5 = 8$).
    \item $\mathbf{6x - 15}$ (car $3 \times 2x - 3 \times 5$).
    \item \textbf{5 cm} (D'après le théorème de Pythagore : $BC^2 = 3^2 + 4^2 = 9 + 16 = 25$, donc $BC = \sqrt{25} = 5$).
    \item $\mathbf{\frac{2}{6}}$ ou $\mathbf{\frac{1}{3}}$ (Les issues favorables sont 5 et 6, soit 2 issues sur 6).
    \item $\mathbf{4,5 \times 10^4}$ (On décale la virgule de 4 rangs vers la gauche).
    \item $\mathbf{x(x + 6)}$ (Le facteur commun est $x$).
    \item \textbf{9 cm} (D'après le théorème de Thalès : $\frac{AD}{AB} = \frac{AE}{AC}$ soit $\frac{2}{6} = \frac{3}{AC}$, donc $AC = \frac{6 \times 3}{2} = 9$).
    \item \textbf{7} (car $f(3) = 3^2 - 2 = 9 - 2 = 7$).
    \item $\mathbf{250\text{ cm}^3}$ (Volume = Aire de la base $\times$ hauteur $= (5 \times 5) \times 10 = 250$).
    \item $\mathbf{x = 3}$ (car $4x = 19 - 7 = 12$, d'où $x = \frac{12}{4} = 3$).
\end{enumerate}

\newpage

% ==========================================
% FICHE 5
% ==========================================
\begin{center}
    \Large\textbf{Entraînement : Automatismes 05}
\end{center}

\noindent \textbf{Partie 1 \hfill 6 points \hfill 20 minutes}

\medskip
\noindent \textbf{Consigne :} Pour chaque question, recopier sur la copie son numéro et la réponse correspondante. Aucune justification n'est demandée.

\bigskip

\noindent
\begin{tabularx}{\textwidth}{|c|X|p{4cm}|}
\hline
\textbf{N$^{\circ}$} & \textbf{Énoncé} & \textbf{Réponse} \\ \hline
\textbf{1} & Écrire sous la forme d'une seule puissance : $\frac{5^7}{5^2}$. & \\ \hline
\textbf{2} & Quel est l'antécédent de $10$ par la fonction linéaire $g(x) = 2x$ ? & \\ \hline
\textbf{3} & Dans une urne contenant 3 boules rouges et 5 boules bleues, on tire une boule au hasard. Quelle est la probabilité de tirer une boule rouge ? & \\ \hline
\textbf{4} & Développer et réduire : $C = (x + 2)(x + 3)$. & \\ \hline
\textbf{5} & Un cylindre de révolution a une base d'aire $20\text{ cm}^2$ et une hauteur de $6\text{ cm}$. Quel est son volume ? & \\ \hline
\textbf{6} & Donner l'écriture décimale de : $3,2 \times 10^{-3}$. & \\ \hline
\textbf{7} & Le triangle $RST$ est tel que $RS = 6\text{ cm}$, $ST = 8\text{ cm}$ et $RT = 10\text{ cm}$. Ce triangle est-il rectangle ? & \\ \hline
\textbf{8} & Calculer et donner le résultat sous forme irréductible : $\frac{4}{3} \times \frac{5}{2}$. & \\ \hline
\textbf{9} & Un automobiliste roule à une vitesse moyenne de $80\text{ km/h}$. Quelle distance parcourt-il en $45\text{ minutes}$ ? & \\ \hline
\textbf{10} & Résoudre l'équation produit nul : $(x - 4)(2x + 6) = 0$. & \\ \hline
\end{tabularx}

\newpage

\subsection*{Correction (pour s'auto-évaluer)}
\begin{enumerate}[label=\textbf{\arabic*.}]
    \item $\mathbf{5^5}$ (car on soustrait les exposants : $7 - 2 = 5$).
    \item \textbf{5} (On cherche $x$ tel que $2x = 10$, donc $x = 5$).
    \item $\mathbf{\frac{3}{8}}$ (Il y a 3 boules rouges sur un total de $3 + 5 = 8$ boules).
    \item $\mathbf{x^2 + 5x + 6}$ (Double distributivité : $x^2 + 3x + 2x + 6$).
    \item $\mathbf{120\text{ cm}^3}$ (Volume = Aire de la base $\times$ hauteur $= 20 \times 6 = 120$).
    \item \textbf{0,0032} (On décale la virgule de 3 rangs vers la droite... attention, vers la gauche pour un exposant négatif !).
    \item \textbf{Oui, il est rectangle en S} (Le plus grand côté est $RT$. $RT^2 = 100$ et $RS^2 + ST^2 = 36 + 64 = 100$. D'après la réciproque du théorème de Pythagore, il est rectangle).
    \item $\mathbf{\frac{10}{3}}$ (car $\frac{4 \times 5}{3 \times 2} = \frac{20}{6} = \frac{10}{3}$).
    \item \textbf{60 km} (45 minutes représentent les $\frac{3}{4}$ d'une heure. $\frac{3}{4} \times 80 = 60$).
    \item $\mathbf{x = 4}$ \textbf{ou} $\mathbf{x = -3}$ (Un produit est nul si l'un de ses facteurs est nul : $x - 4 = 0 \Rightarrow x = 4$ ou $2x + 6 = 0 \Rightarrow 2x = -6 \Rightarrow x = -3$).
\end{enumerate}

\end{document}